\PassOptionsToPackage{table,xcdraw}{xcolor}
\documentclass[11pt, a4paper]{article}

\usepackage[T1]{fontenc}
\usepackage[utf8]{inputenc}
\usepackage{tgpagella}
\usepackage[spanish, es-noshorthands]{babel}
\usepackage{booktabs}
\usepackage{tabularx}
\usepackage{amssymb}
\usepackage[most]{tcolorbox}
\usepackage[margin=1.5cm]{geometry}

\definecolor{ucbblue}{RGB}{0, 51, 102}

\begin{document}

\begin{center}
    \Large\textbf{\textcolor{ucbblue}{Registro Selectivo de Cierre de Lab 3 (Documento Interno Docente)}}[cite: 13]
\end{center}

\begin{tcolorbox}[colback=white, colframe=black, title=\textbf{Objetivo Administrativo}]
    Control y seguimiento focalizado de los grupos que aún adeudan evidencia analítica o física del Lab 3, sin frenar la progresión del curso hacia el Lab 4[cite: 13]. No re-entregar esta guía a los estudiantes[cite: 13].
\end{tcolorbox}

\begin{table}[h!]
    \centering
    \renewcommand{\arraystretch}{1.5}
    \footnotesize
    % La suma de los factores de hsize (0.8 + 1.2 + 1.0) es exactamente 3, igual al número de columnas X.
    % \raggedright elimina los Underfull y Overfull \hbox.
    \begin{tabularx}{\textwidth}{|>{\raggedright\arraybackslash\hsize=0.8\hsize}X|>{\raggedright\arraybackslash\hsize=1.2\hsize}X|>{\raggedright\arraybackslash\hsize=1.0\hsize}X|c|c|}
    \hline
    \rowcolor[HTML]{EFEFEF} 
    \textbf{Identificación del Grupo[cite: 13]} & \textbf{Evidencia Exacta Faltante[cite: 13]} & \textbf{Acción Asignada[cite: 13]} & \textbf{Req. Físico[cite: 13]} & \textbf{Ok[cite: 13]} \\ \hline
    & & & $\square$ Sí \quad $\square$ No & $\square$ \\ \hline
    & & & $\square$ Sí \quad $\square$ No & $\square$ \\ \hline
    & & & $\square$ Sí \quad $\square$ No & $\square$ \\ \hline
    & & & $\square$ Sí \quad $\square$ No & $\square$ \\ \hline
    & & & $\square$ Sí \quad $\square$ No & $\square$ \\ \hline
    & & & $\square$ Sí \quad $\square$ No & $\square$ \\ \hline
    & & & $\square$ Sí \quad $\square$ No & $\square$ \\ \hline
    \end{tabularx}
\end{table}

\end{document}
